\section{Discussion}
\paragraph{Shared operations, task-dependent organization.}
The atlas connects domains through operations that differ in prevalence, timing, and function: verification checks derivations or action preconditions, while revision rejects unoriginal ideas. Matched instructions connect some profiles to task demand. This offers a behavioral connection to the organization of simpler processes in \citet{minsky1986society} and flexible recruitment of shared systems \citep{cole2013hubs}. The evidence concerns textual organization, not discrete mental agents or corresponding brain regions.

\paragraph{Characteristic profiles and useful interventions.}
The transfer matrix separates characteristic profiles from beneficial changes: a math-derived intervention can help science or planning, while amplifying the target's own profile need not help. Diminishing gains at longer cutoffs connect utility to the benefits and costs of computation \citep{shenhav2013control,lieder2020resource}. Natural and forced answering further distinguish producing an answer under a deadline from performance when one is elicited. These effects do not identify an optimal allocation policy or transferred domain knowledge.

\paragraph{Behavioral convergence and its limits.}
Atlas steering and reward learning produce similar mean operation shifts from nearly orthogonal residual vectors. This resembles degeneracy, in which structurally different elements support similar functions \citep{edelman2001degeneracy}, but the eight-label summary could also obscure different processes. RL preserves more verification and improves early forced-answer performance where A2 does not; neither difference establishes mediation. The atlas reveals agreement while the outcome protocols expose distinctions that average profiles miss.

\paragraph{Connections across objectives.}
The math-profile vector lowers paired originality, whereas a math-style instruction raises it while reducing feasibility. Their different effects rule out treating the two as equivalent implementations of mathematical thinking. They do not establish a universal opposition between analysis and creativity; human creative cognition involves cooperation between spontaneous and controlled processes \citep{beaty2015creative}. Cross-task evaluation makes completion and judged quality distinct consequences of reasoning control.

\paragraph{Hypotheses and future work.}
The findings motivate three hypotheses. First, transfer may depend on a target's unmet computational demands rather than source-domain similarity. A selection rule frozen before evaluation should be tested on fresh problem families and instruction paraphrases against own-domain, fixed-vector, random-direction, and stopping baselines. Second, an operation's value may depend on the reasoning stage. Event-triggered control could be compared with fixed schedules and constant steering at matched budgets, measuring candidate generation, error repair, and termination separately. Third, similar average signatures may hide differences in effective verification or candidate diversity. Blinded functional annotation, matched observation windows, and independently trained interventions could distinguish convergence in useful operations from similarity in discourse. These are hypotheses for subsequent studies, not established mechanisms or validated control rules.

\section{Conclusion}
Across seven models, shared operations form reproducible domain signatures in prevalence and timing, with some differences responsive to task demand. Cross-domain steering depends on direction, budget, and objective; a characteristic profile need not be the most useful target. Atlas steering and reward learning converge behaviorally yet differ in early-answer performance. Thought Atlas connects these findings through a common description of reasoning organization and its modification.
